\chapter{Wyniki i dyskusja}\label{ch:wyniki}

\TemplateNote{Rozdział pokazuje tabelę z długimi skrótami (makro \latexcmd{hash} w kolumnie \code{Y}), tabelę ze znacznikami
\latexcmd{cmark} i \latexcmd{xmark} oraz środowisko \code{qa} do zapisu pytań i odpowiedzi.}

\section{Zestawienie wyników}\label{sec:zestawienie}

Tabela~\ref{tab:skroty} zawiera wartości skrótów dla trzech obiektów. Długie ciągi znaków zapisuje się makrem
\latexcmd{hash\{...\}} w kolumnie rozciągliwej, dzięki czemu zawijają się w obrębie komórki, a nie wystają poza margines.

\begin{table}[H]
    \centering
    \caption{Zestawienie sum kontrolnych badanych obiektów}\label{tab:skroty}
    \tablefont
    \begin{tabularx}{\linewidth}{l l Y}
        \toprule
        \hcell{Obiekt} & \hcell{Algorytm} & \hcell{Wartość skrótu} \\
        \midrule
        \multirow{2}{*}{\file{plik_a.bin}}
            & MD5     & \hash{d41d8cd98f00b204e9800998ecf8427e} \\
            & SHA-256 & \hash{e3b0c44298fc1c149afbf4c8996fb92427ae41e4649b934ca495991b7852b855} \\
        \midrule
        \multirow{2}{*}{\file{plik_b.bin}}
            & MD5     & \hash{d41d8cd98f00b204e9800998ecf8427e} \\
            & SHA-256 & \hash{e3b0c44298fc1c149afbf4c8996fb92427ae41e4649b934ca495991b7852b855} \\
        \bottomrule
    \end{tabularx}
\end{table}

Tabela~\ref{tab:wymagania} pokazuje zestawienie spełnienia wymagań ze znacznikami wyboru.

\begin{table}[H]
    \centering
    \caption{Spełnienie wymagań przez badane rozwiązania}\label{tab:wymagania}
    \tablefont
    \begin{tabular}{l c c c}
        \toprule
        \hcell{Wymaganie} & \hcell{Rozwiązanie A} & \hcell{Rozwiązanie B} & \hcell{Rozwiązanie C} \\
        \midrule
        Powtarzalność wyników   & \cmark & \cmark & \xmark \\
        Dokumentacja            & \cmark & \xmark & \xmark \\
        Dostępność narzędzi     & \cmark & \cmark & \cmark \\
        \bottomrule
    \end{tabular}
\end{table}

\section{Odpowiedzi na pytania}\label{sec:pytania}

Pytania i odpowiedzi zapisuje się środowiskiem \code{qa}. Numeracja pytań jest ciągła w całym dokumencie.

\begin{qa}{Czy wyniki są powtarzalne?}
    Tak. Powtórzenie badania w tych samych warunkach dało identyczne wartości skrótów (tabela~\ref{tab:skroty}),
    co potwierdza powtarzalność procedury opisanej w rozdziale~\ref{ch:metodyka}.
\end{qa}

\begin{qa}{Które rozwiązanie spełnia wszystkie wymagania?}
    Spośród zestawionych w tabeli~\ref{tab:wymagania} rozwiązań wszystkie wymagania spełnia rozwiązanie A.
    Odpowiedź można rozwinąć w kilku akapitach, listach lub podsekcjach.
\end{qa}

\section{Dyskusja}\label{sec:dyskusja}

W tej sekcji omów uzyskane wyniki, ich ograniczenia i możliwe źródła błędów. Odwołuj się do rysunków,
tabel i listingów z poprzednich rozdziałów (np. rysunek~\ref{fig:pojedynczy} oraz listing~\ref{lst:kod}),
aby czytelnik mógł łatwo prześledzić tok rozumowania.
